%!TEX root = ../main.tex
%!TEX program = xelatex
% Chapter 3: Words 501 - 750
\chapter{Expanding Expression}
\label{chap:chapter3}

{\small\sffamily\color{mutedtext} Verbs of action, descriptive adjectives, basic concepts, and common social vocabulary.\par}

\vspace{0.4em}
\markboth{Chapter 3: Expanding Expression}{Words 501--750}

\begin{multicols}{2}
\vocentry{501}{Река}{re'ka}{A river}{Эта река опасная: она глубокая и течение очень сильное.}{This river is dangerous: it’s deep and the current is very strong.}
\vocentry{502}{Военный}{va'jennyj}{A military man; military}{Мой папа военный, поэтому мы часто переезжаем.}{My father is a military man; that is why we often move around.}
\vocentry{503}{Мера}{'mera}{A measure, a degree}{Эта мера очень радикальная, но это единственный выход.}{This measure is very radical, but it’s the only way out.}
\vocentry{504}{Страшный}{'strashnyj}{Terrible, scary}{Страшный ураган разрушил всю деревню.}{A terrible hurricane destroyed the whole village.}
\vocentry{505}{Вполне}{vpal'ne}{Quite, fully}{Я не вполне понимаю, что вы имеете в виду.}{I don’t quite understand what you mean.}
\vocentry{506}{Звать}{'zvat’}{To call; to invite}{Я не люблю звать людей по фамилии.}{I don’t like to call people by their last name.}
\vocentry{507}{Произойти}{praizaj'ty}{To take place, to happen (perfective)}{Такое событие могло произойти только в Европе.}{Such an event could take place only in Europe.}
\vocentry{508}{Вперёд}{vpe'rjed}{Forward; in advance}{Мы должны двигаться вперёд и верить в наши идеалы.}{We must move forward and believe in our values.}
\vocentry{509}{Медленный}{'medlennyj}{Slow}{Этот поезд такой медленный! Когда мы прибудем на нашу станцию?}{This train is so slow! When will we arrive at our station?}
\vocentry{510}{Возле}{'vozle}{Beside, by, near}{Возле моей кровати стоит небольшая книжная полка.}{There’s a little book shelf standing beside my bed.}
\vocentry{511}{Никак}{ni'kak}{In no way, by no means}{Эту картину никак нельзя назвать шедевром.}{This painting is in no way a masterpiece.}
\vocentry{512}{Заниматься}{zani'matsa}{To engage in; to study}{Он хочет заниматься научными исследованиями.}{He wants to engage in scientific research.}
\vocentry{513}{Действие}{'dejstvije}{An action}{Некоторые родители контролируют каждое действие своего ребёнка.}{Some parents control every action their child takes.}
\vocentry{514}{Довольно}{da'volnə}{Quite; enough}{Он довольно сдержанный человек.}{He’s quite a reserved man.}
\vocentry{515}{Вещь}{vestch}{A thing; a matter}{Как ты мог взять эту вещь без разрешения?}{How could you take this thing without asking?}
\vocentry{516}{Необходимый}{neobho'dimyj}{Necessary, required}{У вас есть необходимый инструмент, чтобы закончить работу?}{Do you have the necessary tool to finish the work?}
\vocentry{517}{Ход}{hod}{A move; course}{Это был очень хорошо продуманный ход.}{That was a very well-thought move.}
\vocentry{518}{Боль}{bol’}{Pain}{Ужасная боль не давала ей спать по ночам.}{A terrible pain prevented her from sleeping at night.}
\vocentry{519}{Судьба}{syd’'ba}{Destiny, fate}{Ты уверена, что этот мужчина – твоя судьба?}{Are you sure this man is your destiny?}
\vocentry{520}{Причина}{pri'china}{A reason, a cause}{Какова причина вашей ссоры?}{What’s the reason for your quarrel?}
\vocentry{521}{Положить}{pala'zhit’}{To put, to lay down (perfective)}{Не забудь положить мороженое в холодильник.}{Don’t forget to put the ice-cream into the fridge.}
\vocentry{522}{Едва}{jed'va}{Barely}{У него очень плохое произношение, поэтому я едва понимаю, что он говорит.}{He’s got a very bad pronunciation, that’s why I barely understand what he says.}
\vocentry{523}{Черта}{tcher'ta}{A feature, a trait; a line}{Интуитивный интерфейс – это лучшая черта этого приложения.}{Intuitive interface is the best feature of this application.}
\vocentry{524}{Девочка}{'devachka}{A girl}{Девочка была напугана и не могла перестать плакать.}{The girl was scared and couldn’t stop crying.}
\vocentry{525}{Лёгкий}{'ljehkij}{Light, lightweight; easy}{Прими лёгкий душ и расслабься.}{Take a light shower and relax.}
\vocentry{526}{Волос}{'voləs}{A hair}{Официант, у меня волос в тарелке!}{Waiter, there’s a hair in my plate!}
\vocentry{527}{Купить}{ku'pit’}{To buy, to purchase}{Мы не можем позволить себе купить эту машину.}{We can't afford to buy this car.}
\vocentry{528}{Номер}{'nomer}{A number}{Я опять потерял номер твоего телефона.}{I lost your phone number again.}
\vocentry{529}{Основной}{asnav'noj}{Major, basic}{Мы делаем основной упор на развитие экономики.}{We place major emphasis on the development of the economy.}
\vocentry{530}{Широкий}{shi'rokij}{Wide, broad}{Мы предлагаем широкий выбор товаров.}{We offer a wide choice of goods.}
\vocentry{531}{Умереть}{ume'ret’}{To die (perfective)}{Я не готов умереть за политические идеи.}{I’m not ready to die for political ideas.}
\vocentry{532}{Плохо}{'plohə}{Bad, badly}{Ты плохо выглядишь. Что случилось?}{You look bad. What’s happened?}
\vocentry{533}{Глава}{gla'va}{A head (someone in charge, not part of the body); a chapter}{Глава государства должен уметь брать на себя ответственность.}{The head of state must be able to take on responsibility.}
\vocentry{534}{Красивый}{kra'sivyj}{Beautiful}{Красивый закат окрасил море в красный цвет.}{A beautiful sunset painted the sea in red color.}
\vocentry{535}{Серый}{'seryj}{Grey, gray}{Этот серый плащ делает тебя старше. Сними его!}{This grey raincoat makes you look older. Take it off!}
\vocentry{536}{Пить}{'pit’}{To drink}{Я не могу пить такой крепкий кофе.}{I can’t drink such strong coffee.}
\vocentry{537}{Командир}{kaman'dir}{A commander, a commanding officer (military)}{Командир приказал солдатам отступать.}{The commander ordered the soldiers to retreat.}
\vocentry{538}{Обычно}{a'bychnə}{Usually}{Обычно отчёт занимает у меня больше времени.}{The report usually takes me more time.}
\vocentry{539}{Партия}{'partija}{A party (political group); a batch}{Демократическая партия проиграла на прошлых выборах.}{The democratic party lost at the last elections.}
\vocentry{540}{Проблема}{prab'lema}{A problem}{Поверь мне, порванное платье – это не проблема.}{Believe me, a torn dress is not a problem.}
\vocentry{541}{Страх}{strah}{A fear}{Страх потерять работу сделал её очень нервной.}{The fear of losing her job made her very nervous.}
\vocentry{542}{Проходить}{praha'dit’}{To pass, to pass by}{Дорога к дому будет проходить через зелёную лужайку.}{The path to the house will pass through a green lawn.}
\vocentry{543}{Ясно}{'jasnə}{Clearly, distinctly}{В инструкции ясно говорится, что прибор нельзя использовать в воде.}{The instruction clearly states that the device can’t be used in water.}
\vocentry{544}{Снять}{snjat’}{To take off; to take down (perfective)}{Ребёнок попросил маму снять ему обувь.}{The child asked his mom to take off his shoes.}
\vocentry{545}{Бумага}{bu'maga}{Paper}{Эта бумага сделана из переработанной древесины.}{This paper is made of recycled wood.}
\vocentry{546}{Герой}{ge'roj}{A hero; a character}{Мой дедушка – герой войны, и мы им очень гордимся.}{My grandfather is a war hero and we’re very proud of him.}
\vocentry{547}{Пара}{'para}{A couple, a pair}{У меня есть пара идей насчёт новой гостиницы.}{I have a couple ideas about the new hotel.}
\vocentry{548}{Государство}{gasu'darstvə}{A state (government)}{Любое государство должно защищать интересы своих граждан.}{Any state must protect the interests of its citizens.}
\vocentry{549}{Деревня}{de'rjevnja}{A village; the countryside}{Раньше эта деревня была самой большой в районе.}{This village used to be the largest in the area.}
\vocentry{550}{Речь}{'retch}{A speech, a language}{Его речь полна грубых слов и грамматических ошибок.}{His speech is full of rude words and grammar mistakes.}
\vocentry{551}{Начаться}{na'tchatsa}{To begin, to start (perfective)}{Пьеса должна была начаться пятнадцать минут назад.}{The play should have begun fifteen minutes ago.}
\vocentry{552}{Средство}{'sredstvə}{A means}{Мобильный телефон – это самое популярное средство связи сегодня.}{A mobile phone is the most popular means of communication today.}
\vocentry{553}{Положение}{pala'zhenije}{A situation; a location}{Не расстраивайся! Наше положение ещё не худшее.}{Don’t get upset! Our situation is not yet the worst one.}
\vocentry{554}{Связь}{svjaz’}{A link, a connection}{Какова связь между этими двумя ограблениями?}{What’s the link between these two robberies?}
\vocentry{555}{Скоро}{'skorə}{Soon}{Скоро мы узнаем результаты экзамена.}{We’ll know the exam results soon.}
\vocentry{556}{Небольшой}{nebal'shoj}{Small, little}{За домом есть небольшой сад со скамейкой и качелями.}{There’s a small garden behind the house with a bench and a swing.}
\vocentry{557}{Представлять}{predstav'ljat’}{To imagine; to represent; to introduce}{Мне не нужно представлять, как выглядит это здание – я его видел.}{I don’t need to imagine what this building looks like – I’ve seen it.}
\vocentry{558}{Завтра}{'zavtra}{Tomorrow}{Я не думаю, что погода завтра изменится к лучшему.}{I don’t think the weather will change for the better tomorrow.}
\vocentry{559}{Объяснить}{abjas'nit’}{To explain (perfective)}{Невозможно объяснить, что такое любовь.}{It’s impossible to explain what love is.}
\vocentry{560}{Пустой}{pus'toj}{Empty}{Пустой дом выглядел холодным и чужим.}{The empty house looked cold and alien.}
\vocentry{561}{Произнести}{praiznes'ti}{To pronounce, to articulate (perfective)}{Я не могу произнести некоторые русские слова.}{I can’t pronounce some Russian words.}
\vocentry{562}{Человеческий}{thela'vetcheskij}{Human (adj); humane}{Человеческий организм – это сложный механизм.}{A human organism is a complicated mechanism.}
\vocentry{563}{Нравиться}{'nravitsa}{To like (the person who likes takes the dative case)}{Как тебе могут нравиться эти картины?}{How can you like these paintings?}
\vocentry{564}{Однажды}{ad'nazhdy}{One day, once, once upon a time}{Однажды ты встретишь свою судьбу и будешь счастлива.}{One day you’ll meet your destiny and you will be happy.}
\vocentry{565}{Мимо}{'mimə}{Past (preposition)}{Он прошёл мимо нас, как будто никогда не знал нас.}{He went past us as if he never used to know us.}
\vocentry{566}{Иначе}{i'natche}{Otherwise, or else}{Тебе следует уделять больше внимания учёбе, иначе ты не сдашь экзамен.}{You should pay more attention to the studies, otherwise you won’t pass the exam.}
\vocentry{567}{Существовать}{sustchestva'vat’}{To exist}{Жизнь на Земле не может существовать без воды.}{Life on Earth can’t exist without water.}
\vocentry{568}{Класс}{klas}{A class; a classroom; a grade}{Дети, ваш класс самый шумный в школе!}{Kids, your class is the noisiest in the school!}
\vocentry{569}{Удаться}{u'datsa}{To turn out well, to be a success}{Это представление просто не могло удаться – оно было ужасным!}{That performance just couldn’t turn out well – it was horrible!}
\vocentry{570}{Толстый}{'tolstyj}{Fat; thick}{Я и подумать не мог, что он был такой толстый.}{I couldn’t even think he used to be so fat.}
\vocentry{571}{Цель}{tsel’}{A goal, an aim; a purpose}{Наша цель собрать как можно больше информации.}{Our goal is to gather as much information as possible.}
\vocentry{572}{Сквозь}{skvoz’}{Through}{Это окно настолько грязное, что я ничего сквозь него не вижу.}{The window is so dirty that I don’t see anything through it.}
\vocentry{573}{Прийтись}{prij'tis’}{To fall on (a date) (perfective)}{Фестиваль должен прийтись на те же даты.}{The festival must fall on the same date.}
\vocentry{574}{Чистый}{'tchistyj}{Clean; tidy}{Я вижу, ты убиралась дома – пол такой чистый.}{I see you’ve tidied the house – the floor is so clean.}
\vocentry{575}{Знать}{znat’}{To know}{Она была за границей и поэтому не могла знать о том, что случилось.}{She was abroad, so she couldn’t know about what had happened.}
\vocentry{576}{Прежний}{'prezhnij}{Former, previous}{Наш прежний учитель был более дружелюбный.}{Our former teacher was more friendly.}
\vocentry{577}{Профессор}{pra'fesər}{A professor}{Профессор всегда был требователен к своим студентам.}{The professor was always demanding to his students.}
\vocentry{578}{Господин}{gəspa'din}{A gentleman, mister}{Этот уважаемый господин много жертвует на благотворительность.}{This respected gentleman donates a lot to charity.}
\vocentry{579}{Счастье}{'shastje}{Happiness}{Каждый понимает счастье по-разному.}{Every person understands happiness in a different way.}
\vocentry{580}{Худой}{hu'doj}{Thin, skinny}{Мальчик должен есть больше – он слишком худой.}{The boy should eat more – he’s too thin.}
\vocentry{581}{Дух}{duh}{A spirit; a ghost}{Ничто не смогло сломить его сильный дух.}{Nothing could break his strong spirit.}
\vocentry{582}{План}{plan}{A plan}{Он придумал безупречный план.}{He came up with an impeccable plan.}
\vocentry{583}{Чужой}{tchu'zhoj}{Someone else’s; alien}{Это чужой кошелёк, он просто похож на мой.}{It’s someone else’s purse, it just resembles mine.}
\vocentry{584}{Зал}{zal}{A hall; a living room}{Зал был украшен цветами и лентами.}{The hall was decorated with flowers and ribbons.}
\vocentry{585}{Представить}{preds'tavit}{’ - To imagine; to represent; to introduce (perfective)}{Я не могу представить свою жизнь без своих детей.}{I can’t imagine my life without my children.}
\vocentry{586}{Особый}{a'sobyj}{Special, particular}{Я хочу подарить ей особый подарок.}{I want to give her a special present.}
\vocentry{587}{Директор}{di'rectər}{A manager, a director}{Наш директор – очень хороший лидер.}{Our manager is a very good leader.}
\vocentry{588}{Бывший}{'byvshyj}{Ex-, former}{Бывший президент все ещё активно участвует в жизни страны.}{The ex-president still actively participates in the life of the country.}
\vocentry{589}{Память}{'pamjat’}{Memory}{Учёные говорят, что память можно тренировать.}{Scientists say that memory can be trained.}
\vocentry{590}{Близкий}{'blizkij}{Close}{Он мой близкий друг, и я могу доверить ему что угодно.}{He’s my close friend and I can trust him with anything.}
\vocentry{591}{Сей}{sej}{This (archaic, literary, has a humorous connotation if used in everyday speech)}{Ну, что нам говорит сей документ?}{Well, what does this document tell us?}
\vocentry{592}{Результат}{rezul'tat}{A result}{Деньги не проблема – для меня важен результат.}{Money is not a problem; what matters to me is the result.}
\vocentry{593}{Больной}{bal'noj}{Sick, ill; a patient}{Больной ребёнок не спал всю ночь.}{The sick child didn’t sleep the whole night.}
\vocentry{594}{Данный}{'dannyj}{Given, present}{Данный аргумент помог мне доказать, что я прав.}{The given argument helped me to prove that I was right.}
\vocentry{595}{Кстати}{ks'tati}{By the way}{Кстати, не забудь выгулять собаку.}{By the way, don’t forget to walk the dog.}
\vocentry{596}{Назвать}{naz'vat’}{To name, to call (perfective)}{Они ещё не решили, как назвать новую улицу.}{They have not yet decided what to name the new street.}
\vocentry{597}{След}{sled}{Track; a footprint}{Собака потеряла след волка.}{The dog lost the wolf’s track.}
\vocentry{598}{Улыбаться}{uly'batsa}{To smile}{Моя сестра умеет улыбаться, как никто другой.}{My sister can smile like nobody else.}
\vocentry{599}{Бутылка}{bu'tylka}{A bottle}{Стеклянная бутылка упала и разбилась на сотни кусочков.}{A glass bottle fell down and broke into hundreds of pieces.}
\vocentry{600}{Трудно}{'trudnə}{It is difficult, hard}{Трудно совмещать работу и учёбу.}{It is difficult to combine work and studies.}
\vocentry{601}{Условие}{us'lovije}{A condition}{У меня только одно условие – свободный график.}{I have only one condition – a flexible schedule.}
\vocentry{602}{Прежде}{'prezhde}{Before, previously}{После отдыха я чувствую себя гораздо лучше, чем прежде.}{After the rest I feel much better than before.}
\vocentry{603}{Ум}{um}{Mind, intellect}{Кроссворды стимулируют ум.}{Crossword puzzles stimulate the mind.}
\vocentry{604}{Улыбнуться}{ulyb'nutsa}{To smile (perfective)}{Фотограф попросил нас улыбнуться.}{The photographer asked us to smile.}
\vocentry{605}{Процесс}{pra'tses}{A process}{Процесс приготовления этого блюда очень долгий.}{The process of preparation of this dish is very long.}
\vocentry{606}{Картина}{kar'tina}{A picture, a painting}{Эта картина на стене стоит состояние.}{This picture on the wall costs a fortune.}
\vocentry{607}{Вместо}{'vmestə}{Instead of}{Вместо того, чтобы помочь мне, ты даёшь мне бесполезные советы.}{Instead of helping me you are giving me useless advice.}
\vocentry{608}{Старший}{'starshij}{Elder, older, senior}{Её старший брат часто попадает в неприятности.}{Her elder brother often gets into trouble.}
\vocentry{609}{Центр}{tsentr}{Center, centre}{Мы пошли в центр города, чтобы увидеть знаменитый собор.}{We went to the center of the city to see the famous cathedral.}
\vocentry{610}{Подобный}{pa'dobnyj}{Such, similar}{Подобный подход не поможет решить проблему.}{Such an approach won’t help to solve the problem.}
\vocentry{611}{Возможный}{voz'mozhnyj}{Possible}{Возможный союз двух государств может принести много пользы.}{The possible unity of the two states may do a lot of good.}
\vocentry{612}{Около}{'okələ}{Near, around; about}{Твой бумажник на столе, около книг.}{Your wallet is on the table, near the books.}
\vocentry{613}{Смеяться}{sme'jatsa}{To laugh}{Перестань шутить – я больше не могу смеяться.}{Stop joking – I can’t laugh anymore.}
\vocentry{614}{Сто}{sto}{A hundred}{Я сто раз говорила тебе не трогать мои украшения!}{I told you a hundred times not to touch my jewelry!}
\vocentry{615}{Будущее}{'budustcheje}{Future}{Говорят, некоторые люди могут предсказывать будущее.}{They say some people can tell the future.}
\vocentry{616}{Хватать}{hva'tat’}{To be enough (Impersonal, not used in the initial form)}{Тебе хватит этой суммы на неделю?}{Will this sum be enough for you for a week?}
\vocentry{617}{Число}{tchis'lo}{A number; a date}{Ты веришь в то, что семь – это счастливое число?}{Do you believe that seven is a lucky number?}
\vocentry{618}{Всякое}{'fs’akaje}{Things, stuff}{Они будут говорить тебе всякое про меня, но не верь им.}{They're going to tell you things about me but don’t believe them.}
\vocentry{619}{Рубль}{rubl’}{A ruble}{Рубль – это российская валюта.}{The ruble is the Russian currency.}
\vocentry{620}{Почувствовать}{pat'chustvəvat’}{To feel (perfective)}{Это нельзя объяснить, ты должен сам почувствовать это.}{This can’t be explained, you must feel it yourself.}
\vocentry{621}{Принести}{prines'ti}{To bring (perfective)}{Не могла бы ты принести мне печенье из кухни?}{Could you bring me some cookies from the kitchen?}
\vocentry{622}{Вера}{'vera}{A faith, belief; religion}{Её вера в Бога очень сильна.}{Her faith in God is very strong.}
\vocentry{623}{Вовсе}{'vovse}{At all}{Можешь оставить это платье, я не ношу его вовсе.}{You can keep this dress, I don’t wear it at all.}
\vocentry{624}{Удар}{u'dar}{A blow, a strike}{Мощный удар противника нокаутировал боксёра.}{The opponent’s powerful blow knocked the boxer out.}
\vocentry{625}{Телефон}{tele'fon}{A telephone}{По моему мнению, телефон – самое полезное изобретение.}{In my opinion, the phone is the most useful invention.}
\vocentry{626}{Колено}{ka'lenə}{A knee}{Моё колено болит так сильно, что я не могу ходить.}{My knee hurts so much that I can’t walk.}
\vocentry{627}{Согласиться}{sagla'sitsa}{To agree (perfective)}{При всём уважении, я не могу согласиться с вами.}{With due respect I can’t agree with you.}
\vocentry{628}{Коридор}{kari'dor}{A corridor}{Длинный, узкий коридор больницы был пуст.}{The long, narrow corridor of the hospital was empty.}
\vocentry{629}{Мужик}{mu'zhik}{A man (colloquial implying a male with lots of masculinity, able to take on responsibility and do manly deeds)}{Ты можешь положиться на него – он настоящий мужик.}{You can rely on him – he’s a real man.}
\vocentry{630}{Правый}{'pravyj}{Right, right-hand}{Его правый глаз видит лучше, чем левый.}{His right eye can see better than the left one.}
\vocentry{631}{Автор}{'avtər}{Author}{Он автор многих бестселлеров.}{He’s the author of many bestsellers.}
\vocentry{632}{Холодный}{ha'lodnyj}{Cold}{Не пей сок: он слишком холодный.}{Don’t drink the juice: it’s too cold.}
\vocentry{633}{Хватить}{hva'tit’}{To be enough (perfective, unlike ‘хватать’ can be used in the initial form)}{Этого времени должно хватить, чтобы успеть в срок.}{This amount of time should be enough to meet a deadline.}
\vocentry{634}{Многие}{'mnogije}{Many (people)}{Многие говорят, что жизнь стала тяжелее.}{Many are saying that life has become harder.}
\vocentry{635}{Встреча}{'vstrecha}{A meeting, an encounter}{Встреча с этим человеком изменила всю его жизнь.}{The meeting with this man changed all his life.}
\vocentry{636}{Кабинет}{kabi'net}{An office, a private office}{Кабинет начальника на втором этаже.}{The manager’s office is on the second floor.}
\vocentry{637}{Документ}{daku'ment}{A document}{Мы рады, что наконец подписали этот документ.}{We’re happy that we’ve signed this document at last.}
\vocentry{638}{Самолёт}{sama'ljot}{An airplane, a plane}{Самолёт срочно приземлился из-за плохой погоды.}{The airplane landed urgently because of the bad weather.}
\vocentry{639}{Вниз}{vniz}{Down, downwards}{Не смотри вниз, иначе упадёшь!}{Don’t look down or you’ll fall!}
\vocentry{640}{Принимать}{prini'mat’}{To take; to admit}{Ты должен принимать эти витамины два раза в день.}{You should take these vitamins two times a day.}
\vocentry{641}{Игра}{ig'ra}{A game}{Какая твоя любимая компьютерная игра?}{What’s your favorite computer game?}
\vocentry{642}{Рассказ}{ras'kaz}{A story, a narrative; a short story}{Её рассказ о поездке оказался очень увлекательным.}{Her story about the trip turned out to be very exciting.}
\vocentry{643}{Хлеб}{hleb}{Bread}{Она на диете и не ест хлеб.}{She’s on a diet and doesn’t eat any bread.}
\vocentry{644}{Развитие}{raz'vitije}{Development}{Физическое развитие ребёнка зависит от правильного питания.}{Physical development of a child depends on proper nutrition.}
\vocentry{645}{Убить}{u'bit’}{To kill}{Я не уверена, что смогу убить живое существо.}{I’m not sure I’ll be able to kill a living being.}
\vocentry{646}{Родной}{rad'noj}{Native; dear}{Мой родной город много значит для меня.}{My native town means much to me.}
\vocentry{647}{Открытый}{atk'rytyj}{Open; public}{Спасибо за открытый диалог.}{Thank you for the open dialog.}
\vocentry{648}{Менее}{'meneje}{Less}{Этот учебник менее сложный для начинающих.}{This textbook is less difficult for beginners.}
\vocentry{649}{Предложить}{predla'zhit’}{To offer; to suggest; to propose (perfective)}{Я могу предложить вам более доступный вариант.}{I can offer you a more affordable option.}
\vocentry{650}{Жёлтый}{'zholtyj}{Yellow}{В вазе стоит яркий жёлтый подсолнух.}{There’s a bright yellow sunflower in the vase.}
\vocentry{651}{Приходится}{pri'hoditsa}{To have to (impersonal)}{Я живу за городом, поэтому мне приходится рано вставать, чтобы успеть на работу.}{I live in the country, which is why I have to wake up early to be in time for work.}
\vocentry{652}{Выпить}{'vypit’}{To drink; to have a drink (perfective)}{Кто мог выпить весь лимонад?}{Who could drink all the lemonade?}
\vocentry{653}{Крикнуть}{'kriknut’}{To cry out, to scream (perfective)}{Попробуй крикнуть громче, может, они не слышат тебя.}{Try to cry out louder, maybe they don’t hear you.}
\vocentry{654}{Трубка}{'trubka}{A tube; a roll}{Там в ящике лежит металлическая трубка. Передай мне её, пожалуйста.}{There’s a metal tube there in the drawer. Pass it over to me, please.}
\vocentry{655}{Враг}{vrag}{An enemy, a foe}{Лень – мой худший враг.}{Laziness is my worst enemy.}
\vocentry{656}{Показывать}{pa'kazyvat’}{To show}{Мы не хотим никому показывать ребёнка, пока не покрестим его.}{We don’t want to show the baby to anyone before we baptise him.}
\vocentry{657}{Двое}{'dvoje}{Two (in a group together), two of}{Вы двое, идите за мной!}{You two, follow me!}
\vocentry{658}{Доктор}{'doktər}{A doctor (medical)}{Доктор настаивал на операции.}{The doctor insisted on the surgery.}
\vocentry{659}{Ладонь}{la'don’}{A palm (of the hand)}{Цыганка попросила показать ей мою правую ладонь.}{The gypsy woman asked me to show her my right palm.}
\vocentry{660}{Вызвать}{'vyzvat’}{To call, to send for (perfective)}{Мы должны сейчас же вызвать скорую.}{We must call the ambulance at once.}
\vocentry{661}{Спокойно}{spa'kojnə}{Calmly, quietly}{Пожалуйста, объясни всё спокойно и без крика.}{Explain everything calmly and without screaming, please.}
\vocentry{662}{Попросить}{papra'sit’}{To ask for, to ask, to request (perfective)}{Он пришёл, чтобы попросить о помощи.}{He came to ask for help.}
\vocentry{663}{Наука}{na'uka}{Science}{Есть очень много вещей, которые наука не может объяснить.}{There’re lots of things science can’t explain.}
\vocentry{664}{Лейтенант}{lejte'nant}{A lieutenant}{Молодой лейтенант мечтал о большой карьере.}{The young lieutenant dreamed of a great career.}
\vocentry{665}{Служба}{'sluzhba}{Service; employment}{Военная служба сделала его более организованным человеком.}{Military service made him a better organized person.}
\vocentry{666}{Оказываться}{a'kazyvatsa}{To find oneself; to turn out}{Никто не любит оказываться в неловких ситуациях.}{Nobody likes to find themselves in awkward situations.}
\vocentry{667}{Привести}{prives'ti}{To bring (someone with); to lead to (perfective)}{Можно мне привести с собой сестру на вечеринку?}{May I bring my sister to the party with me?}
\vocentry{668}{Сорок}{'sorək}{Forty}{Ему сорок, но он выглядит гораздо старше.}{He’s forty but he looks much older.}
\vocentry{669}{Счёт}{schot}{An account; a bill}{Банковский счёт компании заблокирован.}{The bank account of the company is blocked.}
\vocentry{670}{Возвращаться}{vazvra'stchatsa}{To return, to come back}{Мы не хотим возвращаться в свою старую квартиру.}{We don’t want to return to our old flat.}
\vocentry{671}{Золотой}{zala'toj}{Gold, golden}{Этот золотой медальон принадлежал моей бабушке.}{This gold medallion used to belong to my grandmother.}
\vocentry{672}{Местный}{'mesnyj}{Local}{Этот мужчина – местный гений.}{This man is a local genius.}
\vocentry{673}{Кухня}{'kuhnja}{A kitchen; cuisine}{Наша кухня просторная и светлая.}{Our kitchen is spacious and light.}
\vocentry{674}{Крупный}{'krupnyj}{Large, massive}{Мы получили крупный заказ, так что придётся много работать.}{We got a large order, so we’ll have to work a lot.}
\vocentry{675}{Решение}{re'shenije}{A decision; a solution}{У вас есть два дня, чтобы принять решение.}{You have two days to make a decision.}
\vocentry{676}{Молодая}{mala'daja}{Young (this form is used with feminine nouns)}{Она ещё молодая женщина, но уже многого добилась в жизни.}{She’s still a young woman but has already achieved much in life.}
\vocentry{677}{Тридцать}{'tritsat’}{Thirty}{На улице минус тридцать градусов.}{It’s minus thirty degrees outside.}
\vocentry{678}{Роман}{ra'man}{A novel; a love affair}{Этот роман такой длинный! Я не уверена, что смогу дочитать его до конца.}{This novel is so long! I’m not sure I’ll be able to read it till the end.}
\vocentry{679}{Требовать}{'trebavat’}{To demand, to require; to call for}{Ты не можешь требовать понимания, если сам не умеешь понимать.}{You can’t demand understanding if you yourself can’t understand.}
\vocentry{680}{Компания}{kam'panija}{A company (an enterprise)}{Наша компания быстро растёт.}{Our company is growing fast.}
\vocentry{681}{Частый}{'tchastyj}{Frequent}{Он частый гость в нашем доме.}{He’s a frequent guest in our house.}
\vocentry{682}{Российский}{ras'sijskij}{Russian (when referring to the country)}{Российский рынок активно расширяется.}{The Russian market is expanding actively.}
\vocentry{683}{Рабочий}{ra'botchij}{A worker}{Мой дядя обычный рабочий, но с ним всегда интересно поговорить.}{My uncle is an ordinary worker but it’s always interesting to talk to him.}
\vocentry{684}{Потерять}{pate'rjat’}{To lose}{Я всегда боюсь потерять паспорт, когда путешествую.}{I’m always afraid to lose my passport when travelling.}
\vocentry{685}{Течение}{te'tchenije}{A current; a tendency}{Течение горной реки было быстрым и шумным.}{The current of the mountain river was fast and noisy.}
\vocentry{686}{Синий}{'sinij}{Deep blue}{Я купила синий коврик для ванной.}{I bought a deep blue mat for the bathroom.}
\vocentry{687}{Столько}{'stol’kə}{So much/many}{Я столько хочу тебе рассказать.}{There’s so much I want to tell you.}
\vocentry{688}{Тёплый}{'tjeplyj}{Warm}{Они поженились в тёплый июньский день.}{They got married on a warm June day.}
\vocentry{689}{Метр}{metr}{A meter/metre}{Она упала, когда до финиша оставался всего один метр.}{She fell down when there was only one meter left to the finish line.}
\vocentry{690}{Достать}{das'tat’}{To reach; to take out}{Я не могу достать коробку на верхней полке. Помоги мне, пожалуйста.}{I can’t reach the box on the top shelf. Help me, please.}
\vocentry{691}{Железный}{zhe'leznyj}{Iron (adj)}{Дом окружал высокий железный забор.}{The house was surrounded by a high iron fence.}
\vocentry{692}{Институт}{insti'tut}{An institute; an institution}{Моя сестра поступила в медицинский институт.}{My sister has entered a medical institute.}
\vocentry{693}{Сообщить}{saab'stchit’}{To let know, to communicate}{Я обещаю сообщить тебе, если что-нибудь изменится.}{I promise to let you know if anything changes.}
\vocentry{694}{Интерес}{inte'res}{Interest}{Мой сын проявляет интерес к насекомым.}{My son takes an interest in insects.}
\vocentry{695}{Обычный}{a'bychnyj}{Usual, common, ordinary}{Это был обычный субботний день, но у меня было особое настроение.}{That was a usual Saturday afternoon, but I had a special mood.}
\vocentry{696}{Появляться}{pajav'ljatsa}{To show up; to appear}{Нам нельзя появляться на работе в джинсах.}{We can’t show up at work in jeans.}
\vocentry{697}{Упасть}{u'past’}{To fall, to fall down}{Не спеши – дорога скользкая и ты можешь упасть.}{Don’t be in a hurry – the road is slippery and you can fall.}
\vocentry{698}{Остальной}{astəl’'noj}{Remaining, the rest of}{Мне нравится первый абзац, но остальной текст скучный.}{I like the first paragraph but the remaining text is boring.}
\vocentry{699}{Половина}{pala'vina}{A half}{Сейчас уже половина пятого.}{It’s half past four already.}
\vocentry{700}{Московский}{mas'kovskij}{Moscow (acts like an adjective)}{Я люблю этот старый московский парк.}{I like this old Moscow park.}
\vocentry{701}{Шесть}{shest’}{Six}{Через шесть дней мы едем на море.}{We’re going to the seaside in six days.}
\vocentry{702}{Получиться}{palu'tchitsa}{To be able to, to work out well (impersonal)}{У тебя должно получиться поладить с ними.}{You should be able to get on well with them.}
\vocentry{703}{Качество}{'katchestvə}{Quality; a feature}{Нас интересует качество, а не количество.}{We’re interested in quality, not in quantity.}
\vocentry{704}{Бой}{boj}{A battle, a combat}{Это был страшный бой, и мы потеряли много солдат.}{That was a terrible battle and we lost a lot of soldiers.}
\vocentry{705}{Шея}{'sheja}{A neck}{Я всегда пользуюсь этой мазью, когда у меня болит шея.}{I always use this ointment when my neck hurts.}
\vocentry{706}{Вон}{von}{There (informal)}{Я вижу его. Он вон там.}{I see him. He’s over there.}
\vocentry{707}{Идея}{i'deja}{An idea}{Это просто блестящая идея! Ты гений!}{It’s just a brilliant idea! You’re a genius!}
\vocentry{708}{Видимо}{'vidimə}{Apparently, evidently}{Ты, видимо, очень расстроилась из-за вашей ссоры.}{You are apparently very upset because of your quarrel.}
\vocentry{709}{Достаточно}{das'tatəchnə}{Enough, sufficiently}{Ты уверен, что мы купили достаточно напитков к вечеринке?}{Are you sure that we’ve bought enough drinks for the party?}
\vocentry{710}{Провести}{praves'ti}{To spend (perfective)}{Где вы планируете провести отпуск этим летом?}{Where are you planning to spend your holiday this summer?}
\vocentry{711}{Важно}{'vazhnə}{It is important; importantly}{Важно понимать, что у людей разные вкусы.}{It is important to understand that people have different tastes.}
\vocentry{712}{Трава}{tra'va}{Grass}{Трава перед домом пожелтела от жары.}{The grass in front of the house turned yellow because of the heat.}
\vocentry{713}{Дед}{ded}{A grandfather; an old man (both informal)}{Наш дед любил порядок и был очень строгим.}{Our grandfather liked good order and was very strict.}
\vocentry{714}{Сознание}{saz'nanije}{Consciousness}{Говорят, что духовные практики помогают расширить сознание.}{They say that spiritual practices help to expand consciousness.}
\vocentry{715}{Родитель}{ra'ditel’}{A parent}{Любой родитель переживает, когда дети болеют.}{Any parent is worried when their kids are sick.}
\vocentry{716}{Простить}{pras'tit’}{To forgive (perfective)}{Ты когда-нибудь сможешь меня простить?}{Will you be ever able to forgive me?}
\vocentry{717}{Бить}{'bit’}{To beat}{Нет никаких оправданий, чтобы бить людей.}{There are no excuses to beat people.}
\vocentry{718}{Чай}{tchaj}{Tea}{Утренний чай помогает мне проснуться даже лучше, чем кофе.}{Morning tea helps me to wake up even better than coffee does.}
\vocentry{719}{Поздний}{'poznij}{Late (adj)}{Это был поздний летний вечер со звёздным небом и тёплым ветром.}{That was a late summer evening with a starry sky and a warm wind.}
\vocentry{720}{Кивнуть}{kiv'nut’}{To nod}{Она попросила меня кивнуть, если я согласен на её предложение.}{She asked me to nod if I agreed with her proposal.}
\vocentry{721}{Род}{rod}{Race; generation}{Злодей в фильме хотел уничтожить весь человеческий род.}{The villain in the movie wanted to destroy the whole human race.}
\vocentry{722}{Исчезнуть}{is'tcheznut’}{To disappear, to vanish}{Мы обязательно найдем твою игрушку! Она не могла просто исчезнуть!}{We’ll surely find your toy! It couldn’t just disappear!}
\vocentry{723}{Тонкий}{'tonkij}{Thin}{Утром на озере появился тонкий слой льда.}{A thin layer of ice appeared on the lake in the morning.}
\vocentry{724}{Немецкий}{ne'metskij}{German; the German language}{Это очень надёжный немецкий инструмент.}{It’s a very reliable German tool.}
\vocentry{725}{Звук}{zvuk}{A sound}{Я люблю звук дождя. Он помогает мне расслабиться.}{I like the sound of the rain. It helps me to relax.}
\vocentry{726}{Отдать}{at'dat’}{To give back, to return; to give away}{Девочка попросила подругу отдать её книгу к понедельнику.}{The girl asked her friend to give her book back by Monday.}
\vocentry{727}{Магазин}{maga'zin}{A shop}{Я иду в магазин. Тебе что-нибудь нужно?}{I’m going to the shop. Do you need anything?}
\vocentry{728}{Президент}{prezi'dent}{A president}{Наш президент управляет страной уже восемь лет.}{Our president has been ruling the country for eight years already.}
\vocentry{729}{Поэт}{po'ɛt}{A poet}{Этот молодой поэт уже достаточно хорошо известен.}{This young poet is quite well-known already.}
\vocentry{730}{Спасибо}{spa'sibə}{Thank you (for), thanks for}{Спасибо за полезный совет!}{Thank you for the useful advice!}
\vocentry{731}{Болезнь}{ba'lezn’}{A disease, an illness}{Рак – это ужасная болезнь, но её можно излечить.}{Cancer is a terrible disease but it can be cured.}
\vocentry{732}{Событие}{sa'bytije}{An event}{Рождение ребёнка – это радостное событие для всей семьи.}{The birth of a child is a happy event for the whole family.}
\vocentry{733}{Помочь}{pa'motch}{To help (perfective)}{Не уверен, смогу ли я помочь, но попытаюсь.}{I’m not sure if I’ll be able to help but I’ll try.}
\vocentry{734}{Кожа}{'kozha}{Skin; leather}{Её кожа обгорела на солнце.}{Her skin has burnt in the sun.}
\vocentry{735}{Лист}{list}{A leaf; a sheet}{Жёлтый лист в ручье напоминал парус.}{The yellow leaf in the stream resembled a sail.}
\vocentry{736}{Слать}{slat’}{To send}{Я обещаю слать письма как можно чаще.}{I promise to send letters as often as possible.}
\vocentry{737}{Вспоминать}{vspami'nat’}{To recall, to remember}{Мы любим вспоминать приятные моменты нашей жизни.}{We like to recall the pleasant moments of our life.}
\vocentry{738}{Прекрасный}{prek'rasnyj}{Beautiful, splendid}{Это прекрасный фотоальбом ручной работы.}{It’s a beautiful handmade photo album.}
\vocentry{739}{Слеза}{sle'za}{A tear}{Каждая твоя слеза – трагедия для меня.}{Every tear of yours is a tragedy for me.}
\vocentry{740}{Надежда}{na'dezhda}{Hope}{Как говорят, надежда умирает последней.}{As they say, hope dies last.}
\vocentry{741}{Молча}{'moltcha}{Silently}{Он положил бумаги на стол и молча вышел из комнаты.}{He put the papers on the table and left the room silently.}
\vocentry{742}{Сильно}{'sil’nə}{Strongly, very much}{Он сильно сожалел о том, что порвал с ней.}{He strongly regretted about breaking up with her.}
\vocentry{743}{Верный}{'vernyj}{Faithful, loyal; true}{Этот пёс – мой самый верный друг.}{This dog is my most faithful friend.}
\vocentry{744}{Литература}{litera'tura}{Literature}{Иностранная литература была моим любимым предметом в университете.}{Foreign literature used to be my favorite subject at university.}
\vocentry{745}{Оружие}{a'ruzhije}{Weapon(s), arms}{Он незаконно хранил оружие и был арестован за это.}{He stored weapons illegally and was arrested for it.}
\vocentry{746}{Готовый}{ga'tovyj}{Ready, prepared}{Нам нужен веб-сайт, полностью готовый к работе.}{We need a ready-to-use website.}
\vocentry{747}{Запах}{'zapah}{A smell}{Я ненавижу запах рыбы.}{I hate the smell of fish.}
\vocentry{748}{Неожиданно}{nea'zhidannə}{Unexpectedly}{Зима оказалась неожиданно холодной.}{The winter turned out to be unexpectedly cold.}
\vocentry{749}{Вчера}{vche'ra}{Yesterday}{Вчера мы наконец-то закончили ремонт спальни.}{Yesterday we finally finished the renovation of the bedroom.}
\vocentry{750}{Вздохнуть}{vzdah'nut’}{To sigh; to take a breath (perfective)}{Что заставило тебя так грустно вздохнуть?}{What made you sigh so sadly?}
\end{multicols}
